\subsection{章末確認問題}
\begin{enumerate}
  \item ソフトウェア品質、ソフトウェア製品品質、ソフトウェアプロセス品質の違いを説明せよ。
  \item 「要求に適合しているが、利用者の本当のニーズを満たしていない」場合、品質が高いと言えるか。理由を述べよ。
  \item 誤り、欠陥、故障の違いを、簡単な例を使って説明せよ。
  \item 品質コストにおける予防費用、評価費用、納品前手戻り費用、納品後外部故障費用の違いを説明せよ。
  \item QMS と SQM の関係を説明せよ。
  \item 品質改善で、欠陥を直すだけでなく根本原因分析が必要な理由を述べよ。
  \item SQA が単なるテストではない理由を説明せよ。
  \item 検証と妥当性確認の違いを説明せよ。
  \item 静的解析、動的解析、形式的解析の使い分けを説明せよ。
  \item レビューと監査の違いを説明せよ。
\end{enumerate}

\begin{pointbox}{この資料の主張}
解答の方向性\par
1 は「対象の違い」で説明する。2 は、要求がニーズを正しく表しているかに着目する。3 は「人間の誤り」「成果物中の欠陥」「実行時に観測される故障」の流れで説明する。4 は「品質へ投資する費用」と「品質不足で発生する費用」を区別する。7 は、SQA が品質計画、レビュー、監査、構成管理、測定、V\&V、是正処置を含むことを述べる。8 は、検証が「正しく作っているか」、妥当性確認が「正しいものを作っているか」であると説明できる。
\end{pointbox}
